\section{A sample-dependent posterior ceiling near the common rate}
\label{sec:near-baseline-posterior}

The Fourier construction supplies a low-rate, globally separated
extension of \emph{every} consistently labeled sample. Thus it gives
an actual-posterior comparison, rather than only a representation bound.

\begin{theorem}[Near-baseline critical posterior ceiling]
\label{thm:near-baseline-posterior-ceiling}
There are universal constants $C_0,C_1$ such that the following holds.
Let $S$ be any consistently labeled sample of at most $n\ge1$ Boolean
inputs, and suppose
\[
 L\ge C_0\log(2+d+n+L).
\]
Put
\[
 U_{d,n,L}=V_*+C_1\frac{\log^2(2+d+n+L)}L.
\]
For the original globally margin-conditioned critical posterior,
\[
 \Pi_{N,1}\{g:J_L(g)>U_{d,n,L}\mid S\}\longrightarrow0
 \qquad(N\to\infty).
\]
The dimension, depth, and realized sample are fixed in this limit.
The constants and the rate ceiling are uniform over the sample.
\end{theorem}

\begin{proof}
Let $p_1,\ldots,p_m$ be the distinct negatively labeled sample points,
where $m\le n$. The Boolean function
\[
 g_S(x)=1-2\sum_{j=1}^m\boldsymbol1_{\{x=p_j\}}
\]
agrees with the sample, including any repeated inputs.
A point indicator has normalized Fourier $\ell^1$ norm exactly one,
since
\[
 \boldsymbol1_{\{x=p\}}=2^{-d}\prod_{i=1}^d(1+p_ix_i).
\]
Therefore $A(g_S)\le1+2m\le1+2n$.
Theorem~\ref{thm:controlled-fourier-seed} gives
$J_L(g_S)\le U_{d,n,L}$ after adjusting universal constants.

At fixed $d$, there are finitely many Boolean functions. Their exact
prior rates and the posterior-rate formula are given by
Corollary~\ref{cor:critical-margin-rates}. The denominator's rate is
at most $J_L(g_S)$. If any consistent functions have rate greater
than $U_{d,n,L}$, the finite set of their rates has a strictly positive
gap above this ceiling. Their total posterior mass consequently tends
to zero. If there are none, the assertion is immediate. This argument
places no restriction on mixtures of tied minimizers.
\end{proof}

Combining this ceiling with the universal lower bound gives the limiting
rate interval
\[
 V_*-C/L\le J_L(g)\le
 V_*+C\frac{\log^2(2+d+n+L)}L
\]
for posterior-typical functions, in the stated fixed-domain width limit.
Small remaining differences inside this interval can still determine
learning.

\begin{corollary}[A target-relative obstruction to learning]
\label{cor:critical-excess-obstruction}
For a Boolean target $f$, input distribution $\mathcal D$, and
$0<\epsilon<1$, define
\[
 J_{L,\epsilon}(f;\mathcal D)=
 \inf\{J_L(g):\Pr_{\mathcal D}(g\ne f)<\epsilon\}.
\]
Under the depth hypothesis of
Theorem~\ref{thm:near-baseline-posterior-ceiling}, if
\[
 J_{L,\epsilon}(f;\mathcal D)>U_{d,n,L},
\]
then every sample of at most $n$ inputs labeled by $f$ yields
width-first posterior Gibbs error with lower limit at least $\epsilon$.
Its correct-label probabilities $p_N(x)$ satisfy
\[
 \liminf_{N\to\infty}\Pr_{\mathcal D}\{p_N(x)<2/3\}
 \ge\max\{0,(3\epsilon-1)/2\}.
\]
\end{corollary}

\begin{proof}
The rate ceiling excludes, asymptotically, every function of error
less than $\epsilon$. Every remaining function has error at least
$\epsilon$, proving the Gibbs assertion. If $e_N$ is the displayed
confidence-error probability, Gibbs error is at most
$(1-e_N)/3+e_N$. Rearranging and taking lower limits proves the second
claim.
\end{proof}

In particular, a polynomial sample guarantee would require suitable
weak approximations whose normalized rate excess
$L[J_L(g)-V_*]$ is only $O(\log^2(d+C+L))$, when the sample size is
polynomial in dimension $d$, teacher complexity $C$, and depth $L$,
and the statistical parameters are fixed. This is a necessary
condition, not a proved approximation theorem for arbitrary small
circuits. One should use a fixed error slack when passing from a
Gibbs guarantee to a particular approximant.

The requirement is not vacuous merely because all rates converge to
$V_*$. The uniform all-function upper bound gives normalized excess
$O(d^2+\log^2 L)$, whereas the ceiling for polynomially many examples
has only squared-logarithmic dependence. A negative learning result
would follow from a suitable approximation-excess lower bound for an
explicit small target. A positive learning result must instead control
inaccurate extensions within this much narrower interval. Neither
conclusion follows from the common leading rate alone.
